\cotefile{12}
\scanpage{2}
\mspage{1}
\uline{Esquisse d'une théorie des motifs}.

1. Revue des théories cohomologiques habituelles des schémas projectifs lisses.

\margin{$*$}L'idée de motif, \inserted{qui doit remplacer la cohomologie entière (la théorie des modules de Lef. sur $\unsure{\mathrm Z|\mathbb Z}$ !).} Réserves : a) Isomotifs, b) Motifs $\frac12$ simples c) Base ponctuelle.

2. Définition de la catégorie $M^+(k)$ des motifs effectifs sur $k$.
C'est une catégorie abélienne \inserted{``$\frac12$ simple''} avec $\otimes$. Foncteurs $T_{\mathrm{DR}}^{\deleted{\unsure{*}}}$, $T_{\mathrm{Hdg}}$ (si $k$ de car nulle) $T_B$ (si $k=\mathbb C$) $T^\ell$ (si $\ell$ \deleted{\unsure{premier}} $\inserted{\ne\operatorname{car} k}$)\deleted{\illegible{2}} les foncteurs sont multiplicatifs. Tous ces foncteurs sont fidèles.

3. Graduation dans la catégorie $M^+(k)$. \inserted{Le produit $\otimes$ compatible avec graduation.} Les foncteurs $h^i(X)$. $M^{\unsure{i0|i,0|+}}(X)$.

4. L'équivalence $M^{+1}(k)^\circ \simeq \mathrm{Isab}(k)$ ($X\leadsto\underline{\mathrm{Pic}}^{00}_{X}$, $A\leadsto h^1(\hat A)$).

\deleted{4. Les motifs $Q(-1)$, et la catégorie $M(k)$ des motifs.}

\margin{$*$}Les motifs de poids $\geq1$ comme généralisation des VA. Le caractère ``discret'' est conservé (s'exprimant également par le fait que si $k\unsure{\longrightarrow|\lhook\joinrel\longrightarrow} k'$ est une extension de corps algébriquement clos, alors $M^{\unsure{e}}(k)\unsure{-|\longrightarrow}M^e(k')$ est pleinement fidèle).

\margin{$*$}Caractère continu\ldots, Exemples d'énoncé suggéré par la théorie des \deleted{\unsure{schémas}} abéliens : caractérisation des motifs \uline{lisses} sur $S$ \deleted{\unsure{régulier (point|réguliers (point} \unsure{exclu)}}\inserted{normal}.

5. Le motif $Q(\unsure{+1|-1})$ (anti-Tate). La catégorie $M(k)$ des motifs (isomotifs semi-simples sur $k$ !). Maintenant il y a des \uline{Hom} internes, en plus \unsure{des|du} $\otimes$ et de la graduation.

Si $C$ est une catégorie pseudo-abélienne avec $\otimes$, alors les foncteurs compatibles avec $\otimes$ de $M(k)$ dans $C$ correspondent aux foncteurs $T:V(k)^\circ\longrightarrow C$, satisfaisant Künneth, et tels que \deleted{\unsure{l'objet}} $T(Q(-1))=\operatorname{Im} T(\pi)$ ($\pi$ le projecteur dans $h(\unsure{\mathrm P|\mathbb P}^1_k)$ \deleted{\unsure{degré de}}, d'image $h^2(\mathrm P^1_k)=Q(-1)$) \uline{soit} \uline{inversible}. On montre qu'un tel $T$ commute automatiquement à la formation du \uline{Hom}. La catégorie des motifs $M(k)$ est donc caractérisée, avec sa structure $\otimes$, \deleted{et l'objet $Q(\unsure{-1})$} à équivalence près par propriété universelle.
% archive sheet number: 1

